% SPDX-License-Identifier: CC-BY-4.0
% Copyright (C) 2026 Fabian Januszewski, Abhijit (the Langlandscape team)
% GENERATED by tools/extract_statements.py from mathematics/chapters/ -- do not edit by hand.

\begin{quoted}{Theorem}{2.1}{31}{thm:s01-mass}{Eichler mass formula; \texttt{if:s01-mass}}{chapters/s01/2-mass.tex}
\begin{enumerate}[label=(\alph*)]
\item $\displaystyle\sum_{[x]\in\XB{\hat\cO^\times}}\frac1{|\Gamma_x|}=\frac{1}{24}\prod_{\ell\mid D}(\ell-1)$
  (full-unit convention), equivalently $\sum 1/w_x=\varphi(D)/12$.
\item (level change) For open $U'\le U\le\hat\cO^\times$: $\sum_{X_B(U')}1/|\Gamma_{x'}(U')|=[U:U']\sum_{X_B(U)}1/|\Gamma_x(U)|$.
  With $w$ in place of $|\Gamma|$ the same holds iff $-1\in U'$ or $-1\notin U$; if $-1\in U\setminus U'$ then
  $2\sum_{X_B(U')}1/w_{x'}=[U:U']\sum_{X_B(U)}1/w_x$ (\lean{Langlandscape.Quat.massW\_level\_change\_two}).
\item $[\hat\cO^\times:\Uzero{N}]=\psi(N)=N\prod_{\ell\mid N}(1+\ell^{-1})$ and $[\hat\cO^\times:\Uone{N}]=\psi(N)\varphi(N)=N^2\prod_{\ell\mid N}(1-\ell^{-2})$.
\item Hence $\sum_{\XB{\Uzero{N}}}1/w_x=\varphi(D)\psi(N)/12$, and for $p\ge5$, $p\nmid D$, $s\ge1$, $(M,pD)=1$:
\[|S_s|=|\XB{\Uone{Mp^s}}|=\frac{\varphi(D)}{24}\,\psi(Mp^s)\,\varphi(Mp^s).\]
\end{enumerate}
Status: (a) \stcited{} \lit{voight\_2021\_quaternion-algebras}{Thm.~25.1.1, p.~415 (with $w_J=\#\cO_L(J)^\times/\{\pm1\}$,
hence $1/12$)}; \lit{eichler\_1955\_zur-zahlentheorie-der-quaternionen-algebren}{\S4 ``Klassenzahl und Maß'', p.~136ff.};
in Lean it is the registered hypothesis \pathtt{Hyp.EichlerMass} (\Tpath{Hyp/EichlerMass.lean}, $1/24$, full units), proved
for $D=2,11$ in tests. (b) \stformal{} \Tpath{P2/Mass.lean} \pathtt{mass\_orbit\_count}; the $w$-variant \stformal{} \pathtt{Quat.massW\_level\_change},
\pathtt{Quat.massW\_level\_change\_two} (orbit counting). (c) first half \stformal{} \pathtt{gl2\_borel\_index}, \pathtt{mass\_level\_change}; second half
\stproved{} below (count $|\GL_2(\Z/N)|/(N\varphi(N))$). (d) first half \stpmc{} from (a)--(c) (proof below; cf.\ \lit{voight\_2021\_quaternion-algebras}{Thm.~25.3.18,
p.~423}), \stpmc{} in Lean as \pathtt{eichlerOrder\_mass} (under \pathtt{Hyp.EichlerMass}); second half \stpmc{} from (a)--(c)
and Lemma~1.7(b) (all $\Gamma_x=1$), and \stformal{} under \pathtt{Hyp.EichlerMass} as \pathtt{Quat.card\_XB\_Utower}.
\end{quoted}

\begin{quoted}{Theorem}{6.1}{48}{thm:s01-JLclassical}{classical JL, Brandt form; \pathtt{if:s01-targets}; seed Thm.~3.1}{chapters/s01/6-targets.tex}
Let $(M,D)=1$; for the $\Gamma_0$ statement let $k\ge2$ be even. Choosing $\Qp\hookrightarrow\C$ (or working with the import's global splitting over a number
field), there is an isomorphism of modules over the algebra generated by the $T_\ell$, $\ell\nmid DM$,
\[\Brandt{\Uzero{M}}{V_k}\otimes\C\ \cong\ S_k(\Gamma_0(DM))^{D\text{-new}}\ \oplus\ \begin{cases}\C\cdot\mathbf 1&k=2,\\0&k>2,\end{cases}\]
where $\mathbf 1$ spans the functions factoring through $\Nrd$ (one-dimensional since $\Q_{>0}\backslash\hat\Q^\times/\Zhat^\times=1$),
and $H^1_{\rm par}(\Gamma_0(DM),V_k\otimes\C)\cong S_k\oplus\overline{S_k}$. At $\Gamma_1$-level, for $(N,D)=1$ and $k\ge2$ of any parity: for $k\ge3$,
$\Brandt{\Uone{N}}{V_k}\otimes\C\cong S_k(\Gamma_1(N)\cap\Gamma_0(D))^{D\text{-new}}$, and for $k=2$ the same after dividing
by the functions factoring through $\Nrd$, compatibly with $T_\ell,S_\ell$ for $\ell\nmid ND$.
\stcited{} \lit{jacquet-langlands\_1970\_automorphic-forms-on-gl2}{Thm.~16.1 (LNM~114, p.~494; p.~250 in the retypeset copy in \pathtt{literature/}) (representation-theoretic; JL call
their proof a sketch)}; \lit{buzzard\_2004\_p-adic-families-of-automorphic-forms\_preprint}{Thm.~2, p.~13 ($\Gamma_1$ form,
attributed to JL, Shimizu, Arthur)}; \lit{hijikata-pizer-shemanske\_1980\_basis-problem-for-modular-forms-announcement}{\S0,
p.~280 (basis problem, all $N$)}; \lit{shimizu\_1965\_zeta-functions-of-quaternion-algebras}{(scan; not text-verified)};
\lit{pizer\_1980\_an-algorithm-for-computing-modular-forms-on-gamma0-n\_scan}{(scan; not text-verified; algorithmic side)};
survey \lit{voight\_2021\_quaternion-algebras}{Remark~41.5.13, p.~778 (historical account of the correspondence; the
statement itself is not formulated there)}.
\par\smallskip\noindent{\scriptsize\itshape [Status paragraph following the statement in the chapter:]}\par
\noindent\emph{Lean status.} \lean{Langlandscape.Hyp.JLClassical} \stformal{} as a \emph{schematic interface} (S01-L13): only the case $k=2$
with trivial coefficients, the Eisenstein line $\C$ with $T_\ell=\ell+1$, and an abstract Hecke module in place of
$S_k(\Gamma_0(DM))^{D\text{-new}}$. It is a statement about the true space only once that space is substituted for the abstract datum, which cannot be checked
in Lean here; no later result may be advertised as ``under classical JL'' without this caveat.
\end{quoted}

\begin{quoted}{Theorem}{6.3}{49}{thm:s01-lambdaJL}{$\Lambda$-adic Hecke JL}{chapters/s01/6-targets.tex}
Let $\T^{\ord}(DM)$ be Hida's ordinary $\Lambda$-adic Hecke algebra of tame level $\Gamma_1(M)\cap\Gamma_0(D)$ and
$\T^{\ord}_B(M)$ the one acting on $\varprojlim_se\,\Brandt{U_s}{\cOE}$ (or on its homological dual). There is a surjection
$\T^{\ord}(DM)\twoheadrightarrow\T^{\ord}_B(M)$ with $T_\ell\mapsto T_\ell$ ($\ell\nmid DMp$), $U_\ell\mapsto U_\ell$ ($\ell\mid Mp$),
classical $\langle d\rangle\mapsto$ seed $\langle d\rangle^{-1}$ for $d$ prime to $Mp$ (the guard $\gcd(d,Mp)=1$, stated in \lean{Langlandscape.Hyp.LambdaJL} as \texttt{Nat.Coprime d (M*p)}, S04-L17a; the diamonds are defined only for such $d$; the seed's adelic diamond, right translation by $\mathrm{diag}(1,d)$; Remark~4.5, Lemma~5.5), and $U_\ell\mapsto+\Pi_\ell$ ($\ell\mid D$), $\Pi_\ell$ acting by right translation; Hsieh writes $-$, which is not imported (stage 05, \pathtt{rem:s05-hsieh}),
\provisional{F6/D3: sign fixed to $+\Pi_\ell$, Hsieh's sign declared not imported; before: $U_\ell\mapsto\pm\Pi_\ell$ ($\ell\mid D$; sign convention-dependent); the Hsieh citation below read ``$\mathrm{JL}(U_q)=-U_q$ for $q\mid N^-$'' without the remark that it is not imported.}
whose kernel cuts out the
$D$-new quotient: $\T^{\ord,D\text{-new}}\cong\T^{\ord}_B$.
\stcited{} \lit{longo-vigni\_2009\_quaternion-algebras-heegner-points-and-the-arithmetic-of-hida-families\_arxiv}{Prop.~6.1,
p.~27 (iso of $\tilde\Lambda$-algebras $\T^{\ord}_\infty\simeq\mathbb B^{\ord}_\infty$, with $T(n,n)=\langle n\rangle$ normalised as in
their [25])}; \lit{hsieh\_2017\_hida-families-and-p-adic-triple-product-l-functions\_arxiv}{\S4.5, p.~34 (surjection
$\mathrm{JL}:\T(N,\mathbb I)\to\T^D(N,\mathbb I)$, with his sign $\mathrm{JL}(U_q)=-U_q$ for $q\mid N^-$, not imported)}.
\par\smallskip\noindent{\scriptsize\itshape [Status paragraph following the statement in the chapter:]}\par
\noindent\emph{Lean status.} \lean{Langlandscape.Hyp.LambdaJL} \stformal{} as a \emph{schematic interface} (S01-L13): two conjuncts, the surjection
$\varphi$ with the generator matching (including classical $\langle d\rangle\mapsto$ seed $\langle d\rangle^{-1}$ and $U_\ell\mapsto+\Pi_\ell$ for $\ell\mid D$; the diamond field of the Lean hypothesis is corrected to the inverse diamond on the Lean side, S04) and $\ker\varphi=$ an abstract ideal cutting out the $D$-new
quotient. The $\GL_2$ side (algebra, generators, ideal) is abstract data, so the hypothesis is $\Lambda$-adic JL only after instantiation by the true objects.
\end{quoted}
